%%%%%%%%%%%%%%%%%%%%%%%%%%%%%%%%%%%%%%%%%%%%%%%%%%%%%%%%%%%%%%%%%%%%%%%%
%                   GLOSARIO                                           %
%%%%%%%%%%%%%%%%%%%%%%%%%%%%%%%%%%%%%%%%%%%%%%%%%%%%%%%%%%%%%%%%%%%%%%%%
\newglossaryentry{ether}{
name=Ether,
description={Ether (ETH) es la criptomoneda nativa de la blockchain de Ethereum. Se utiliza para pagar tarifas de transacción (gas) y ejecutar contratos inteligentes dentro de la red. También actúa como un activo digital que puede ser intercambiado o utilizado en aplicaciones descentralizadas (DApps)}
}\newglossaryentry{wei}{
name=wei,
description={El wei es la unidad más pequeña de Ether (ETH), donde 1 ETH equivale a \(10^{18}\) wei. Es la unidad de medida utilizada en las transacciones de la blockchain de Ethereum}
}
\newglossaryentry{metamask}{
name=MetaMask,
description={MetaMask es una extensión de navegador que permite a los usuarios interactuar con aplicaciones descentralizadas (DApps) en la red Ethereum, facilitando la gestión de cuentas y la firma de transacciones}
}
\newglossaryentry{oraculo}{
name=Oráculo,
description={Un oráculo en blockchain es un servicio que proporciona datos externos a los contratos inteligentes, permitiéndoles interactuar con información del mundo real, como precios, clima o eventos deportivos}
}
\newglossaryentry{web3}{
name=Web3,
description={Web3 es la evolución de la web que utiliza tecnología blockchain para descentralizar aplicaciones y servicios, permitiendo mayor control, privacidad y propiedad de los datos por parte de los usuarios}
}
\newglossaryentry{blockchain}{
name=blockchain,
description={Blockchain es una tecnología de registro distribuido que almacena datos en bloques encadenados de forma segura, inmutable y descentralizada. Se usa principalmente para transacciones digitales, como criptomonedas, pero también tiene aplicaciones en contratos inteligentes, logística y más}
}
\newglossaryentry{addresstype}{
name=Tipo Address,
description={El tipo address representa direcciones de la red}
}
\newglossaryentry{euribor}{
name=Euribor,
description={Es el tipo de interés promedio al que los principales bancos de la zona euro se prestan dinero entre sí en el mercado interbancario. Se calcula diariamente y sirve como referencia para numerosos productos financieros, especialmente para hipotecas a tipo variable en Europa}
}
\newglossaryentry{smartcontract}{
name=Smart Contract,
plural=Smart Contracts,
description={Programa o código autoejecutable que se ejecuta en una blockchain. Define reglas y condiciones preestablecidas y, cuando se cumplen, ejecuta automáticamente acciones sin necesidad de intermediarios. Se usa en transacciones, acuerdos digitales y aplicaciones descentralizadas.}
}
\newglossaryentry{script}{
name=script,
plural=scripts,
description={Conjunto de instrucciones o código escrito en un lenguaje de programación o de scripting que se ejecuta automáticamente para realizar tareas específicas}
}
\newglossaryentry{javascript}{
name=JavaScript,
description={JavaScript es un lenguaje de programación interpretado, ligero y versátil, utilizado principalmente para crear contenido web interactivo}
}
\newglossaryentry{sha256}{
name=SHA-256,
description={Secure Hash Algorithm 256 es una función hash criptográfica que genera un valor único de 256 bits (32 bytes) a partir de una entrada}
}
\newglossaryentry{bitcoin}{
name=Bitcoin,
description={Red blockchain original basada únicamente en transacciones y la que mayor valor ha conservado}
}
\newglossaryentry{hash}{
name=hash,
description={Un hash es un valor único generado por una función hash a partir de un conjunto de datos de entrada. Se utiliza para representar esos datos de manera compacta y segura}
}
\newglossaryentry{hashing}{
name=hashing,
description={Proceso de realizar un hash}
}
\newglossaryentry{dapp}{
name=dApp,
description={Decentralized App o aplicación descentralizada, es una aplicación que funciona sobre una red descentralizada de blockchain en lugar de un servidor central. Utiliza contratos inteligentes para gestionar su lógica y datos}
}
\newglossaryentry{appstore}{
name=App Store,
description={Tienda digital de aplicaciones}
}

\newglossaryentry{ethereum}{
    name = Ethereum,
    description= {Red blockchain cuya moneda principal es el ETH y es capaz de ejecutar \glspl{smartcontract}}
}
\newglossaryentry{eth}{
    name= ETH,
    description= {Ether es la criptomoneda de la red \gls{ethereum}}
    }
\newglossaryentry{gas}{
    name= gas,
    description={Unidad que mide el trabajo computacional necesario para ejecutar operaciones en Ethereum. La comisión se paga en ETH y depende del gas consumido y del precio por unidad de gas}
}
\newglossaryentry{defi}{
    name= DeFi,
    description={Decentralized Finance: Sistema financiero en un entorno descentralizado, como puede ser una red blockchain}
}
\newglossaryentry{UI}{
    name= UI,
    description={User Interface: Interfaz de usuario}
}
\newglossaryentry{framework}{
    name= Framework,
    description={Conjunto de herramientas, bibliotecas y convenciones diseñadas para facilitar el desarrollo de aplicaciones, proporcionando una estructura base que simplifica y organiza el código}
}
\newglossaryentry{solidity}{
    name= Solidity,
    description={Lenguaje de programación utilizado para programar contratos inteligentes}
    }
\newglossaryentry{nft}{
name= NFT,
description={Token no fungible: activo digital distinguible de otros tokens por un identificador propio. No es intercambiable de forma equivalente con cualquier otro token. El estándar ERC-721 permite registrar su titularidad y transferirlo; su relación con un activo externo depende del sistema que lo utilice~\cite{erc-721}}
}
\newglossaryentry{token}{
name= Token,
description={Activo digital emitido en una red blockchain que representa un valor, un derecho o una utilidad específica dentro de una plataforma}
}

\newglossaryentry{pool}{
name= Pool,
plural= Pool,
description={Grupo de fondos disponibles para prestar en una plataforma DeFi.}
}
\newglossaryentry{wallet}{
name= Wallet,
plural= Wallets,
description={Software o hardware que te permite almacenar, gestionar y firmar 
transacciones con tus criptomonedas sin depender de un intermediario como un banco.}
}
\newglossaryentry{aave}{
    name= AAVE,
    description={Protocolo de finanzas descentralizadas (DeFi) de código abierto y no custodio que permite a los usuarios depositar criptoactivos para obtener rendimientos o tomarlos prestados mediante contratos inteligentes sin intermediarios}
    }
\newglossaryentry{uniswap}{
    name= Uniswap,
    description={
        Bolsa de intercambio descentralizada basada en Ethereum que automatiza el intercambio de criptomonedas mediante fondos de liquidez, sin intermediarios ni libros de órdenes tradicionales}
    }
\newglossaryentry{criptomoneda}{
    name={criptomoneda},
    description={Activo digital que utiliza técnicas criptográficas y una red blockchain para registrar y verificar transacciones sin necesidad de una autoridad central. Puede emplearse como medio de intercambio}
}
\newglossaryentry{pop-up}{
    name={Pop-up},
    description={Ventana emergente que aparece sobre la interfaz de una aplicación o página web para mostrar información o solicitar una acción del usuario, como la autorización de una transacción}
}
\newglossaryentry{diagramauml}{
    name={Diagrama-UML},
    description={Un diagrama UML (Unified Modeling Language) es una representación gráfica utilizada para modelar, visualizar, documentar y diseñar sistemas de software}
}
\newglossaryentry{secuenciauml}{
    name={Diagrama de secuencia UML},
    description={
        Diagrama que muestra cómo interactúan los objetos de un sistema 
en un orden temporal, representando el intercambio de mensajes entre ellos}
}
\newglossaryentry{login}{
    name={Log-In},
    description={
        El login (o inicio de sesión) es el proceso de autenticación mediante el cual un 
        usuario se identifica para acceder a un sistema, aplicación, sitio web o dispositivo 
        informático de forma segura}
}




%%%%%%%%%%%%%%%%%%%%%%%%%%%%%%%%%%%%%%%%%%%%%%%%%%%%%%%%%%%%%%%%%%%%%%%%
%                    ABREVIATURAS                                      %
%%%%%%%%%%%%%%%%%%%%%%%%%%%%%%%%%%%%%%%%%%%%%%%%%%%%%%%%%%%%%%%%%%%%%%%%
\newacronym{ods}{ODS}{Objetivos de Desarrollo Sostenible}
\newacronym{op}{OP}{Objetivo principal}
\newacronym{os}{OS}{Objetivo secundario}
\newacronym{pow}{PoW}{Proof of Work, Prueba de trabajo}
\newacronym{pos}{PoS}{Proof of stake, Prueba de participación}
\newacronym{evm}{EVM}{Ethereum Virtual Machine, Máquina Virtual de Ethereum}
\newacronym{dni}{DNI}{Documento Nacional de Identidad}
\newacronym{ocde}{OCDE}{Organización para la Cooperación y el Desarrollo Económico}
\newacronym{amear}{AMEAR}{Acotado, Medible, Específico, Alcanzable, Relevante}
\newacronym[
    longplural={organizaciones no gubernamentales},
    shortplural={ONG}
    ]{ong}{ONG}{organización no gubernamental}
\newacronym[
    longplural={proyectos de fin de grado}
    shortplural={PFG}
    ]{pfg}{PFG}{proyecto de fin de grado}
\newacronym[
    description={S.P.E.C.I.A.L es la sigla usada para los atributos de
        Fuerza (\textbf{S}trenght), Percepción (\textbf{P}erception),
        Resistencia (\textbf{E}ndurance), Carisma, (\textbf{C}harisma),
        Inteligencia, (\textbf{I}ntelligence), Agilidad,
        (\textbf{A}gility), y Suerte (\textbf{L}uck)
    },
    ]{special}{SPECIAL}{\textit{Strenght, Perception, Endurance,
        Charisma, Intelligence, Agility \& Luck}
    }
